\section{Introduction}
\label{sec:introduction}
Lobbying networks that undermine public trust in climate science and delay effective policy action are extensively documented by investigative journalists and civil society organisations, including dedicated platforms such as DeSmog~\cite{desmog} and InfluenceMap's LobbyMap~\cite{lobbymap}, which together publish nearly 2000 profiles of people and organisations who lobby against climate action and policy.
Despite the breadth of such work, the resulting information remains siloed across heterogeneous platforms and is published primarily in long-form prose, organised around organisations' and individuals' profiles.
As a consequence, even basic questions, such as who funds a given think tank, which staff previously worked for fossil fuel companies, or which board members sit across multiple front groups, require tedious manual cross-referencing across sites with incompatible structures.
Investigative journalists, NGOs, and watchdog organisations rely on this kind of manual analysis, which does not scale.
Representing this domain as a knowledge graph makes such relational knowledge explicit and queryable at scale.

In this paper, we present LACK (Lobbying Agents, Climate, and Knowledge), a knowledge graph of persons and collectives involved in lobbying against climate science and policy, built from over 2k source documents and guided by a purpose-built lightweight ontology.
We make three contributions.
First, we present the LACK resource: a purpose-built vocabulary covering the key relation types in the climate lobbying domain, and a knowledge graph of typed relations between persons and organisations, extracted from authoritative sources, linked to Wikidata and DBpedia, and published as Linked Open Data.
Second, we demonstrate the analytical value of LACK through case studies that surface connections spanning multiple independent sources, and we assess the resource both for its coverage of the competency questions behind its design and through a survey of domain experts in investigative journalism and climate accountability.
Third, drawing on both the construction process and this expert engagement, we contribute an analysis of the opportunities and open challenges of building such resources automatically in this domain, including temporal change, funding traceability, shifting positions, long-tail entities, and geographic skew.

The remainder of this paper is structured as follows.
Section~\ref{sec:related} situates LACK in the related literature.
Section~\ref{sec:lack} describes how LACK is built and what it contains, covering the ontology, the extraction process, and the resulting resource.
Section~\ref{sec:evaluation} reports the case studies and the expert evaluation.
Section~\ref{sec:discussion} discusses the opportunities and open challenges of building and using such a resource in practice, and concludes with directions for future work.